\chapter{先学会安全地走一步}\label{sec:descent}
\section{平方和下降，到底还缺什么}
令 $E(x)=\sum_i x_i^2$。平均一组 $S$，其均值为 $a$，则
\[
E(x)-E(T_Sx)=\sum_{i\in S}x_i^2-pa^2
=\sum_{i\in S}(x_i-a)^2.
\]
右边非负，且除非选中的值全相等，否则严格为正。这是“平均让数更接近”的精确版本。

要从这里得到有限终止，还缺两个条件。第一，中间数最好都是整数，这样下降量不可能无穷趋近于零。第二，别走到上一讲的共同余数陷阱里。我们将把这两个条件和一种位置储备一起保持：\textbf{有两个不同数值，各出现至少 $p$ 次。}

这种数值叫“重值”。比如五平均中 $a$ 出现六次，它是重值；出现四次则不是。重值不是说 $a$ 的绝对值大。

\section{留在外面的两个数，为什么能保护这一步}
假如未选位置里还留着一个奇数和一个偶数，那么不管选中的位置怎样变化，输出都不可能全同奇偶。模任意素数也是同一个道理。因此，与其直接检查新的均值，不如事先在外面保留一对不同余位置。

什么时候这会紧张？只有绝大多数位置都属于同一个余数类、少数例外很容易被一次全拿走的时候。
\InsertProof{03_descent_tools-unit1}
例如五平均二十元中，若十八个位置是偶数、两处是奇数，那么不要把两处奇数同时拿走；保留一处，就保留了奇偶差异。对每个危险素数留一个例外位置，总共最多留 $d_p(n)$ 个。
\InsertProof{lem:safety}
这里还给了几个更省位置的选择。最常用的是 $(a^{p-1},x)$：留下的一份 $a$ 会迫使任何假想的共同余数等于 $a$，而平均方程再迫使 $x$ 也等于这个余数，矛盾。正是这一份重复值，省掉了另找见证的麻烦。

\section{为什么要保留两个重值}
一次平均会产生 $p$ 份相等的新值。如果另外还保留着一个旧的 $p$ 重值，而且新旧两值不同，就又回到“两个重值”的状态，可以继续下一轮。

注意，新均值若恰好等于保留的旧值，就只剩一个重值。这种情况叫“碰撞”。它不是说原配置不可解，只是当前的储备没有保持住。下一讲会通过碰撞等式重新选一组。

初始化时，在分析中先考虑 $pX$。所有初值都被 $p$ 整除，所以第一次任取 $p$ 项的均值是整数。选出位置序列后，在原 $X$ 上执行同一序列即可。后面的整数尺度固定下来，不在下降途中随意改变。

一个小技巧反复有用：对 $p+1$ 个候选数，删掉一个后剩 $p$ 个。如果所有删法都得到同一个和，比较 $S-x_i$ 与 $S-x_j$ 就知 $x_i=x_j$。因此，多留一个候选位置便有机会避开指定均值。
\InsertProof{lem:initialize}

\section{整除要求交给一个零和定理}
安全性安排好以后，还要使所选和被 $p$ 整除。只看模 $p$ 的余数，这就是找 $p$ 项零和。Erdős--Ginzburg--Ziv 定理在此处的结论很直观：$2p-1$ 个候选永远够；少一个时，失败的形状也能描述。
\InsertProof{03_descent_tools-unit4}
\InsertProof{lem:egz}

第一次读这个证明，可以把它拆成两幅图景。前半段用 $z_i^{p-1}$ 表示“第 $i$ 项选不选”：它在零处为零，其余地方为一。第一方程迫使选项数是 $p$，第二方程迫使所选余数之和是零。有限域求和证明这样的选择确实存在。

后半段处理少一个候选的边界。把一些位置两两配对，每对可以选左边或右边。这提供一串二选一的自由度。每增加一对，可能的余数和至少增加一种；只要一开始比普通情况多出一种，就能覆盖全部 $p$ 个余数。唯一躲开这个过程的，是两类各 $p-1$ 份的形状。

到此，我们有了三件工具：留下余数见证保证安全；重复值保证下一轮还有材料；零和定理保证平均仍为整数。下一讲只负责把它们安排成一个覆盖所有状态的选择方法。
